\section{Introduction}

When evaluating implementations for blockchain, virtually any program containing several parties who share products or services and depend on a central authority to handle such transactions may be examined. The transaction to be managed in this specific use case is the rental of a car. As a consequence, the owner of a vehicle can provide a rented car, and the clients can use it for some time in return for payment. There are several car-sharing companies out there, such as Car2go or DriveNow, all of whom rely on a centralized back-end system that authorizes rentals and manages transactions and payments. However, the case of private car sharing, where someone possessing a vehicle will offer rental services to the public, is not unique. Companies like Drivy provide a platform where the car owner can offer their car to potential renters. Unfortunately, providing this kind of service to users brings many security flaws; first of all is the overall stability of service. Using DDoS attacks on the centralized server would lead to the collapse of the whole service. In comparison, when data are distributed and processed over the whole network of nodes like blockchain, it is practically impossible to tamper the service \cite{Nakamoto,Astarita2019}.

The idea of car-sharing is not unique to the car industry, but the usage of blockchain technologies to monitor car ownership, document vehicle utilization, and issue insurance costs, and other transactions is, on the contrary, very recent \cite{Bossauer2019UsingBI,Tsoniotis2019,Valastin2019}. Unlike today's approaches, blockchain-based car-sharing will be both more straightforward and more available to a wider variety of people \cite{Dorri2017}. The intermediary will be substituted by a distributed ledger technologies (DLTs), returning 99\% of the total rental earnings to the owner of the car.

Why should we use a blockchain over traditional systems? That is a very good and valid question. Let us look at the 10,000 feet car sharing example. Let's take Bob here: he has to earn money, so he needs to go to work, quite a common scenario. He could go in a retro way: ride a horse to work or go by taxi, public transport or ride a bike. The unfortunate reality is most likely he will take a car. Let's assume, for a moment, that he is going to use something like a car-sharing this time. So he is going by car from his home to his workplace, which is a few kilometers or miles away. If he wants to use any modern car-sharing, then he has most likely to select a suitable date and time upfront in the calendar that gets saved on a server of the car-sharing company. They also have access to his credit card, which will be charged at some point, probably once a month. So he pays, and he is given access to the car. The car can be unlocked with his phone or a special card, and he can go from A to B. That is pretty straightforward. 

How would that work on a blockchain? One thing we might know already is that everyone has access to the same data. It is a shared ledger after all, and it is public so everybody can at least read. The car knows if Bob paid something, the car knows from smart contract maybe for which period Bob paid upfront, and the car knows the public key of Bob or a wallet. Bob can open the car using his private key, and then Bob can drive. How does that exact same scenario work in a traditional client-server system with a car? The car has access to the Internet, most likely with a SIM card in some sort of data stick or something. Then there is a central server from the car-sharing company, and they have an API, and the car needs to fire a request against an API at a central database. That is again very straightforward. There are a few drawbacks that also mean that data is stored centralized on a server. Even if it is stored in a cloud or across different data centers, it is still stored in such a way that there is a master key to access it. The users have little or almost no control over their own data. It is basically a monopoly running there, and all in all is pretty intransparent what is going on.

Nevertheless, it has been like this for years and pretty much all services we know from Facebook, Twitter to Windows updates, or a favorite weather API, it all is a request against the server containing some data. How will this same thing work on a blockchain? The car still has access to the Internet, and it needs some Internet. However, instead of connecting to a server, it connects to many different peers downloading a database locally. The car is not just downloading the blockchain, it is verifying on-the-fly if the blocks are all valid. This is done using cryptography, which is basically maths on steroids. It is part of what is called consensus. What is consensus actually? Take chess for example, maybe the historical game between Kasparov and Deep Blue. We might or might not know the rules, but even though the game is already long, we can still see at any round if both participants obeyed by the rules of chess. If we see mid-game that someone moved an impossible way, we would know that he/she cheated, and that is really all there is to consensus. Consensus on the blockchain would mean that every participant would run through the whole game again, checking if someone was cheating. 

Back to our car-sharing example, there are a few advantages over traditional client-server methods. First of all, it is all key-less. What it means is that because there is a private-public key pair, there does not need to be any sign up or user ID or whatever we would have in a traditional system. Our wallet would take care of authorizing access to the car, and the car would know it because it downloaded the blockchain's latest state. The second big bonus is a peer to peer payment. We could basically pay the car directly, and the car would know that we paid, so we have access and can use it for a certain amount of time. The third big plus, it is all transparent. The blockchain is an immutable ledger. Every transaction can be traced back to its origin. Last but not least, because every item, everything, has access to the same information, IoT (Internet of Things) can perfectly work with it. The car would magically know that it has to open its locks or even pay for its own gas eventually. There might be many cars, but instead of traditional client-server architecture, we are moving the whole database around and give everyone the same data. That boils down to having the right key, and the key is different for every person. We call it private keys, and they generate and control addresses and so on, but that is a different story. 

In addition, traditional data servers provide us CRUD interface, so we are able to create, read, update, and delete data. This suggests that some people might change users' data without further notification of the users. On the contrary, blockchain as immutable public ledger allows us
only to create and read data. If we want to update some information in blockchain, we have to notify every single user in the network that we have changed data. As a matter of course, it is also impossible to delete data from the blockchain. It can be concluded that blockchain, as an unchangeable method, increases transparency and ensures data protection \cite{xu2016blockchain}.

This article is organized as follows: In Section II, we are describing the design and architecture of our blockchain car-sharing solution with a focus on the tokenized architecture. Section III shows the evaluation and results of the implemented verification of the design. The related works are decomposed in Section IV, and finally, we discuss the results, conclusions, and future work in Section V.
